\chapter{结论}
\label{chap:concl}

本工作面向独居老人安全监护场景，设计并实现了``智护家''多模态健康监护系统，通过融合视频、音频、生理数据、用药记录四种模态，对老人健康状况进行三级风险评估。以下从主要工作总结、核心发现和未来展望三个方面进行总结。

\section{主要工作总结}

本工作历经三个迭代版本，主要完成以下工作：

\begin{enumerate}
    \item \textbf{四模态融合架构设计与实现}：构建了基于4层Cross-Modal Attention的多模态融合网络，各模态通过独立投影层对齐到统一特征空间（768/768/256/128维$\to$512维），经4层跨模态注意力交互后加权融合，输出三级风险分类。模型参数量8.51M，支持注意力权重可视化。

    \item \textbf{真实数据驱动}：摒弃早期版本的纯合成数据，引入ESC-50音频数据集（11类$\to$3风险等级）和Pexels视频数据集（112段，daily/struggle/falls三类），结合半合成生理与用药数据，构建了900样本的多模态数据集（训练720/验证90/测试90）。

    \item \textbf{视频模态真实化}：将视频特征从零向量升级为基于MediaPipe Tasks API PoseLandmarker提取的真实人体骨架运动特征。通过33个关键点的轨迹、躯干角度、质心速度等物理量，经确定性投影生成768维特征向量，特征非零率从7.6\%提升至100\%。

    \item \textbf{端到端应用集成}：开发了基于Streamlit的交互式Web应用，支持视频/音频文件上传、生理数据表单输入、用药记录录入，实时输出风险等级、置信度、模态权重和注意力矩阵热力图，并集成缺失数据默认值处理机制。
\end{enumerate}

\section{核心发现}

\subsection{跨模态注意力优于简单拼接}

消融实验A组表明，基于Cross-Modal Attention的Full模型达到100\%测试准确率，而基于MLP拼接的Lite模型仅为94.44\%。Lite模型在``中风险$\to$高风险''方向出现4次误判、``高风险$\to$中风险''方向出现1次误判，说明跨模态注意力能够有效捕获模态间的交互信息，对风险等级边界附近的判别能力更强。

\subsection{四模态均对决策有贡献}

消融实验B/C组的结果如表\ref{tab:conclusion-ablation}所示：

\begin{table}[H]
    \centering
    \caption{模态贡献消融结果汇总}
    \label{tab:conclusion-ablation}
    \begin{tabular}{lccc}
        \toprule
        掩蔽模态 & 测试准确率 & 准确率下降 & 重要性排序 \\
        \midrule
        完整模型（无掩蔽） & 100.00\% & --- & --- \\
        掩蔽视频 & 98.89\% & 1.11\% & 3 \\
        掩蔽音频 & 98.89\% & 1.11\% & 3 \\
        掩蔽生理数据 & 67.78\% & 32.22\% & 1 \\
        掩蔽用药记录 & 97.78\% & 2.22\% & 2 \\
        \bottomrule
    \end{tabular}
\end{table}

四种模态掩蔽后准确率均下降，验证了``四模态融合''设计的合理性。其中生理数据贡献最大（32.22\%），是区分风险等级的核心信号；用药记录次之（2.22\%）；视频与音频贡献相当（各1.11\%），主要作为辅助验证信号。

\subsection{真实视频特征对模型训练有效}

迭代2（视频零向量）测试准确率为94.74\%，迭代3（真实视频特征）提升至100\%。更重要的是，迭代2中视频模态贡献为0\%（掩蔽视频不影响准确率），而迭代3中视频模态贡献达到1.11\%，证明真实骨架特征使视频模态真正参与了融合决策。

\subsection{端到端推理可靠}

应用层端到端测试中，6组真实ESC-50音频样本全部预测正确，且置信度随风险等级递增（低风险91\%、中风险95\%、高风险97\%），符合临床直觉，验证了从原始多模态输入到最终风险评估的完整链路可靠性。

\section{局限性回顾}

本工作的主要局限包括：（1）数据规模有限（900样本），100\%准确率可能反映测试集规模而非模型泛化能力；（2）视频特征基于骨架关键点而非端到端学习，对姿态遮挡敏感；（3）生理与用药数据为半合成，标签与风险等级存在强相关性；（4）模型在CPU上训练，未充分利用GPU加速；（5）各模态全局权重设为均匀分布，未实现自适应权重学习。这些局限在第\ref{chap:discuss}章中已详细讨论。

\section{未来展望}

基于本工作的基础，未来可从以下方向深入：

\begin{enumerate}
    \item \textbf{数据扩充与真实化}：引入更多真实跌倒视频数据集（如URFD、FDD），采集真实居家场景的生理与用药数据，提升数据多样性与泛化能力。

    \item \textbf{端到端视频编码}：采用VideoMAE等预训练视频编码器替代骨架关键点提取，直接从原始视频帧学习时空特征，提升对遮挡和弱光的鲁棒性。

    \item \textbf{自适应模态权重}：将固定均匀权重改为可学习的模态权重网络，根据输入数据的可靠性和信息量动态调整各模态贡献。

    \item \textbf{时序建模}：引入时序卷积或循环结构，对连续多天的健康数据进行趋势分析，实现风险预警而非单次评估。

    \item \textbf{临床验证}：与合作养老机构开展试点部署，收集真实使用反馈，评估系统在实际环境中的有效性与误报率。
\end{enumerate}

综上所述，本工作构建了一个从数据采集、特征提取、多模态融合到应用展示的完整系统，验证了跨模态注意力机制在老年人健康监护场景中的有效性，并通过严格的消融实验和端到端测试提供了可信的评估证据。尽管存在数据规模等局限，系统架构和方法论为后续研究和实际部署奠定了基础。
